\chapter{Research Methodology and Evaluation Design}
\label{cha:methodology}

A successful demonstration shows that an investigator can solve a
particular incident. Establishing its usefulness requires comparison
with simpler approaches and tests of where it fails. A convincing answer
may change when a record is missing or the endpoint is healthy. This
chapter defines the research approach, comparison methods, independent
cases and measures needed to assess diagnosis and remediation under
controlled conditions.

\section{Research Approach}
\label{sec:methodology:approach}

The thesis follows an engineering research approach: a software artefact
is developed to address an operational problem and assessed through
empirical observations. The ACM SIGSOFT \emph{Empirical Standards for
Software Engineering} describe engineering research in terms of a
justified artefact and its conceptual and empirical evaluation. Their
benchmarking guidance emphasises an explicit workload, meaningful measures
and repeatable comparison \parencite[Engineering Research and
Benchmarking]{sigsoft2026empirical}. Here, the artefact is the investigation
and remediation system. The empirical method combines a controlled
diagnostic comparison with live tests of the operational workflow.

Known faults support development and refinement. Reserved cases test the
fixed approaches without using their answers for development. The
assessment distinguishes the ability to operate the system from the
usefulness of its conclusions.

Functional checks verify the connectors, access controls and operator
workflow. Comparative diagnosis tests assess correctness and evidence
support. Live remediation tests check recovery, approval and rollback.
Passing the functional checks is a prerequisite for interpreting the
comparisons, rather than evidence of diagnostic effectiveness.

Table~\ref{tab:methodology:rq} connects each research question to its
experimental assessment.

\begin{table}[htbp]
\centering\small
\caption{Research questions and their experimental assessment.}
\label{tab:methodology:rq}
\begin{tabularx}{\textwidth}{@{}>{\raggedright\arraybackslash}p{0.07\textwidth}>{\raggedright\arraybackslash}X>{\raggedright\arraybackslash}X@{}}
\toprule
RQ & Assessment & Evidence collected \\
\midrule
RQ1 & Controlled requests with known relationships and competing traffic. & Source availability, supported path links, missing links and unjustified request matches. \\
\addlinespace
RQ2 & Matched investigations with selected security fields and summaries removed or added. & Changes in diagnosis, evidence support and input size, with independently checked feature calculations. \\
\addlinespace
RQ3 & Repeated comparison on reserved fault, healthy and incomplete-evidence cases. & Component and cause accuracy, evidence support, abstention, failures and resource use. \\
\addlinespace
RQ4 & Controlled plan review, approval, execution, recovery and rollback tests. & Accepted and rejected actions, actual state changes, endpoint recovery and rollback verification. \\
\bottomrule
\end{tabularx}
\end{table}

\section{Comparison Methods and Controls}
\label{sec:methodology:comparison}

The proposed comparison contains three approaches: a rule-based procedure,
a single investigating agent and a collaborating agent configuration.
They address two different questions. Comparing the procedure with the
single agent examines the value of model-directed investigation beyond
explicit troubleshooting rules. Comparing the two agent configurations
examines whether collaboration adds value under the same resource limits.
Neither comparison assumes that the more complex approach performs better.

\subsection{Rule-Based Reference Procedure}

The baseline is an executable troubleshooting procedure whose next check
and conclusions follow explicit conditions. It uses the same telemetry
tools as the agents and returns the same diagnostic fields. For example,
a blocking security event is relevant only if its target and time match
the investigation and the request was intended to be permitted. A nearby
denial alone does not establish the cause. Missing observations lead to
further checks or an insufficient-evidence result, rather than a default
diagnosis.

The rules are developed from operational knowledge and the known
development cases. A Norlys engineer reviews their coverage, stopping
conditions and treatment of ambiguity before the final cases are used.
The procedure may branch according to observations. It does not require
every investigation to begin at the firewall or follow one linear path.
Its rules and any mappings from diagnoses to repair recommendations are
fixed alongside the agent configurations.

This baseline represents one documented procedure, not an average
engineer. It provides a reproducible comparison without a human-participant
experiment. Consequently, the measured execution time cannot establish
how much time an engineer would save, and differences in query order are
not themselves errors.

\subsection{Single and Collaborating Investigators}

The single investigator can select tools, revise its interpretation and
request further evidence. It is allowed multiple model calls. Restricting
it to one response would compare different opportunities to investigate
as well as different agent organisations.

The collaborating configuration distributes investigative responsibilities
and combines their findings into one final result. Both configurations
use the same model version, underlying tool capabilities, evidence
representation and output schema. The single investigator receives the
same substantive diagnostic instructions and reference material, although
the collaborating configuration distributes them between roles. The role
division and communication mechanism are design choices. The experimental
contrast concerns the implemented organisation as a whole, rather than
agent count isolated from every accompanying change.

The comparison uses shared limits on total model tokens, tool calls and
elapsed time per investigation. All agents, coordination messages,
retries and final-answer generation consume that total allowance. A token
allowance applies across model input and output, including reported
reasoning usage without counting it twice. Unreported usage is recorded
as a measurement limitation. The rule-based procedure shares the tool
and time limits but has no model-token consumption. A method may finish
early and is not required to spend its allowance.

The input contains the same endpoint, time interval, symptom description
and system context. No method receives another method's answer. Each run
starts with empty incident memory. Development examples, if included,
are identical in informational content and exclude the reserved cases.
Different query choices remain part of the methods being evaluated.

\subsection{Model Selection and Pilot}

Model selection precedes the final comparison. Candidate eligibility is
checked against N2 and N5: permitted deployment and processing locations,
tool use, structured output, version identification and available quota.
The eligible candidates are tested on the same development cases with
the same single-investigator workflow. The pilot records diagnostic and
evidence scores using the rubric below, together with invalid tool calls,
latency and token use. Selection considers supported diagnoses first,
then the observed resource trade-offs. A successful connection alone is
not a selection result.

The pilot also establishes feasible investigation limits and estimates
the cost of repeated evaluation. Limits must allow each agent
configuration to complete its intended workflow on development cases.
The chosen model, settings and limits are fixed before reserved-case
execution. Model selection is therefore not another variable in the
main agent comparison. Its candidates, observations and selection
rationale belong with the system design.

\section{Cases and Independent Reference Answers}
\label{sec:methodology:cases}

An experimental case consists of a defined system condition, a request
workload, an observation interval and an independent reference record.
A run is one approach's attempt to investigate that case. Repeating a
run measures variation in the investigator. It does not create a new
fault or an independent application environment.

Cases are selected deliberately to exercise the diagnostic challenges,
rather than sampled in proportion to the frequency of production
incidents. The principal fault cases contain one injected cause so that
its effect can be checked. Similar symptoms should be represented by
different causes where feasible, preventing an HTTP status alone from
determining the answer. Table~\ref{tab:methodology:cases} defines the
required distinctions in the case set.

\begin{table}[htbp]
\centering\small
\caption{Case groups and the behaviour they test. These are selection
criteria, not a record of completed injections.}
\label{tab:methodology:cases}
\begin{tabularx}{\textwidth}{@{}>{\raggedright\arraybackslash}p{0.23\textwidth}>{\raggedright\arraybackslash}X@{}}
\toprule
Case group & Diagnostic purpose \\
\midrule
Application and platform faults & Distinguish an application failure from an unavailable or incorrectly configured hosting component. \\
\addlinespace
Security-control faults & Identify a control that blocks a legitimate request and explain the mismatch with intended policy. \\
\addlinespace
Healthy controls & Establish expected behaviour even when the investigation begins with a reported failure or irrelevant error. \\
\addlinespace
Incomplete evidence & Test whether removing a decisive source changes the conclusion appropriately without creating false certainty. \\
\addlinespace
Ambiguous correlation & Test competing requests and nearby events for which timing and addresses do not identify one request uniquely. \\
\bottomrule
\end{tabularx}
\end{table}

The selection is constrained by the observation and change boundaries in
Section~\ref{sec:requirements:context}. Shared production components can
supply evidence but cannot be injection targets. A security-control case
requires an authorised, isolated control affecting the test workload.
Missing permissions or unavailable fault mechanisms reduce demonstrated
coverage. They do not justify replacing a live control failure with
invented log records and reporting it as equivalent. Coverage is recorded
against the Norlys fault catalogue \parencite[AC-7]{norlys_operations_agents}.

For every case, the reference record identifies the injected condition,
faulty component, expected request behaviour, actual observed effect and
reset procedure. It also defines acceptable diagnostic specificity and
the observations needed to distinguish the cause. The injected condition
is checked against resource state and independent request observations.
Issuing an injection command without verifying its effect is insufficient
to establish ground truth.

Before diagnostic execution, the case custodian checks whether the retained
evidence can support that reference at the stated level of specificity.
Cases that cannot support it belong to the incomplete-evidence assessment,
with the limitation recorded. They are not silently discarded or moved
between groups after an approach fails. This separates a failure to use
available evidence from an inability to observe the cause.

Known development cases remain available for debugging. Norlys coworkers
prepare the reserved cases and retain their exact conditions and reference
answers until every compared output has been recorded. Fault families
and permitted targets may be known in advance. An unseen case can thus
test a new instance of a familiar failure mechanism without demonstrating
generalisation to an entirely new mechanism.

Test identifiers carry no cause information. Explicit scenario names,
injection records and answer keys are excluded from all diagnostic
inputs. Ordinary operational evidence, including error messages and
relevant configuration, remains available. A realistic symptom that
reveals the cause is valid evidence. A test label supplied solely because
the evaluator knows the answer is not.

\thesisTODO{Complete the case allocation with Norlys: authorised injection
targets, distinct cases per fault family, healthy controls, evidence
variants and eligible repairs. Name the case custodian and reviewer.
Retain exact reserved fault conditions outside the developer's materials.}

\section{Experimental Procedure}
\label{sec:methodology:procedure}

\subsection{Evidence Collection and Repeated Diagnosis}

Live connector checks first establish that each source supplies the
required evidence for the selected test workload. This includes separate
checks for empty results, missing coverage, access denial and query
failure. Controlled requests record their send times, responses and
available identifiers. Cases with overlapping requests test whether
correlation preserves ambiguity. Functional checks also exercise the
interface, distinct viewing and approval permissions, history retrieval
and access to the controls while the test API is unavailable.

For the diagnostic comparison, each case's scoped telemetry and relevant
configuration observations are retained as a fixed evidence collection.
The collection is captured after a predefined ingestion waiting period,
using the same selection and masking rules for every case. It includes
the bounded operational observations, not only records selected because
they support the reference answer. Source gaps and collection failures
are retained explicitly.

Diagnostic tools query this fixed collection so that the available
evidence cannot improve merely because one method runs later. All methods
can access the same observations, although they may retrieve different
subsets. Requesting information outside the collection returns an
explicit unavailable result. Live tests separately verify the connectors
and repair workflow. This separation permits a controlled comparison
of investigation while preserving a test of interaction with the actual
environment. Diagnosis time on retained evidence does not include live
log-ingestion delay.

The final diagnostic procedure is:

\begin{enumerate}
\item Verify the case condition and capture the evidence and reference
record. Keep the reference separate from the diagnostic tools.
\item Run each fixed approach on the same case collection. Interleave
approaches and randomise their order within repetitions, retaining the
schedule and random seed. No investigation changes the target state.
\item Stop at a submitted final result or the first resource limit.
Retain the complete tool and model-call record, final output and reason
for stopping. Do not provide corrective human hints.
\item Record all final diagnoses and proposed repair plans before opening
the reserved reference answers for scoring. Apply the same procedure to
healthy and evidence-variant cases, retaining their distinct group labels
in the analysis.
\end{enumerate}

Each agent configuration has the same number of fresh runs per case.
The deterministic baseline is checked for repeatability, but duplicate
answers are not counted as additional independent diagnostic evidence.
Failures after a valid run begins, including timeouts and malformed
outputs, remain outcomes. A case whose injection never took effect is
an invalid setup, recorded separately before scoring. A later correction
to a reference answer must apply to every method's outputs for that case.

\thesisTODO{Insert the pilot-supported parameter table before final
evaluation: case and repetition counts, token/tool/time limits, retry
policy, ingestion wait, workload rate and concurrency, and recovery
observation duration. State the planned run total and estimated cost.}

\subsection{Security Evidence and Robustness}

RQ2 requires more than listing fields used by the final design. A focused
ablation compares three evidence views on the security cases. An ablation
changes a defined part of the input to assess its contribution. The first
view provides individual events with their basic time and target fields.
The second adds the available security decision and policy context, such
as action, rule and intended permission. The third adds selected summaries
calculated from those same records, such as denied connections to the
target within the incident interval.

The single-investigator configuration is held fixed for this comparison,
including its model, prompt and limits. This avoids changing evidence
and agent organisation together. Field groups, summary definitions and
eligible cases are fixed using development data. Equivalent views cannot
be reconstructed through another tool. An independent calculation checks
each summary's filters, denominator and unit. A fraction of denied
connections is not scored as a fraction of failed HTTP requests.

The analysis compares diagnostic correctness, evidence support and input
size across these views. A benefit from summaries may reflect easier
access to existing information rather than new information. Results are
therefore interpreted as effects of the selected representation in these
cases, not universal feature importance.

Robustness checks use paired variants of a case, for example withholding
one source or adding competing records from the scoped workload. The
reference states whether the remaining evidence still supports a diagnosis
or only a narrower conclusion. Removing any field does not automatically
make abstention correct. Variants remain linked to their original case
and are not treated as new independent faults.

\section{Outcome Measures and Analysis}
\label{sec:methodology:measures}

The diagnostic measures implement F2--F4 and the distinctions established
in Section~\ref{sec:state-of-the-art:evaluation}. Every run returns one
primary conclusion. Listing the correct cause among several incompatible
possibilities does not receive full diagnostic credit. Equivalent wording
is accepted when it describes the reference mechanism at the required
level of specificity.

\subsection{Correctness, Evidence and Uncertainty}

On fault cases, component correctness and cause correctness are recorded
separately as binary outcomes. A component is correct when it matches
the reference target at the specified granularity. A cause is correct
when it explains the injected condition rather than only its symptom.
Joint diagnostic correctness requires both. Thus, naming the application
without identifying its misconfiguration can receive localisation credit
without being a correct diagnosis.

Evidence is assessed separately using
Table~\ref{tab:methodology:evidence}. References are first checked for
existence. Their relevance and support for the decisive claims then
require review against the retained observations. The evidence assessment
does not require one particular record when several records support the
same conclusion.

\begin{table}[htbp]
\centering\small
\caption{Evidence-support rubric applied to the primary conclusion.}
\label{tab:methodology:evidence}
\begin{tabularx}{\textwidth}{@{}>{\raggedright\arraybackslash}p{0.22\textwidth}>{\raggedright\arraybackslash}X@{}}
\toprule
Assessment & Interpretation \\
\midrule
Supported & Decisive claims refer to relevant observations that justify the conclusion and its stated uncertainty. Contradictory observations are addressed. \\
\addlinespace
Partly supported & Relevant observations exist, but a decisive causal link or asserted detail is unsupported. \\
\addlinespace
Unsupported & No usable support is provided, references are fabricated or irrelevant, or the conclusion contradicts the cited observations. \\
\bottomrule
\end{tabularx}
\end{table}

The primary outcome is an evidence-supported correct diagnosis: joint
correctness together with the supported evidence rating. Joint accuracy
without this evidence condition is also reported, making correct guesses
visible. For the company-level correct/partially correct/incorrect account
in AC-7, a jointly correct answer is correct, a correct component with an
incomplete cause is partially correct, and a wrong target or contradictory
cause is incorrect. Abstentions and operational failures remain separate
categories and receive no diagnostic-success credit.

Healthy cases have their own denominator. Their false-diagnosis rate is
the number of runs asserting a fault divided by all validly started healthy
runs. Supported healthy conclusions, insufficient-evidence outcomes and
operational failures are reported separately so that a method that never
answers cannot appear useful merely by avoiding false diagnoses.
On reduced-evidence cases, the assessment records appropriate qualified
conclusions or abstentions and unsupported fault or healthy assertions.

For RQ1, the reviewer checks each claimed path link against the request
record and available identifiers. The results distinguish supported links,
omitted observable links and claims stronger than the observations permit.
An ambiguous connector-to-application connection is not a verified match
to one HTTP request. Where no unique reference match can be established,
the evaluation assesses the uncertainty statement rather than inventing
an exact match for scoring.

After the diagnostic runs, I score the outputs against the reference
records. A Norlys reviewer independently assesses a predefined sample
covering every approach and case group, with method names hidden. The
sample is selected before scores are examined. Original ratings,
agreement and resolved disagreements are retained. Cases referred for
additional review are reported separately from this agreement sample.
This review concerns the reference and evidence rubric, not similarity
to another model's preferred explanation.

\thesisTODO{Specify the independently reviewed sample size and selection
seed with the reviewer. Calibrate the scoring rubric on development
outputs before scoring reserved cases.}

\subsection{Consistency, Resource Use and Comparative Analysis}

For method $m$ and fault case $i$, let $s_{mir}$ equal one when run $r$
produces an evidence-supported correct diagnosis and zero otherwise.
With $R_i$ scheduled repetitions, the case success fraction is
\begin{equation}
\label{eq:methodology:case-success}
p_{mi}=\frac{1}{R_i}\sum_{r=1}^{R_i}s_{mir}.
\end{equation}
Timeouts and invalid outputs contribute zero. The primary aggregate is
the mean of these fractions across the $N$ distinct fault cases,
$\bar p_m=N^{-1}\sum_{i=1}^{N}p_{mi}$, giving each case equal weight.
Healthy controls and deliberately reduced-evidence variants are reported
separately. Each table states both the number of cases and the number of
runs. For the deterministic baseline, a case contributes its single
diagnostic outcome.

Comparison is paired by case. For example,
$\Delta_i=p_{\mathrm{multi},i}-p_{\mathrm{single},i}$ shows whether
collaboration helps or harms on the same observations. The analysis
reports these differences, their mean and the number of cases favouring
each approach. The baseline comparison follows the same principle.
Results are also separated by fault family so that numerous application
cases cannot conceal weak coverage elsewhere.

Consistency is reported through the per-case success fractions and the
frequency of each primary component/cause conclusion across repetitions.
A stable wrong answer remains a failure. With a small, deliberately
selected case set, the analysis is primarily descriptive. Repeated calls
are not treated as independent incidents for significance tests or as a
population estimate of production accuracy.

Resource accounting records input and output tokens, model and tool calls,
elapsed diagnosis time and failures, including retries. Coordination
usage belongs to the collaborating configuration's total. Token cost
estimates identify the applicable prices and date, with cached usage
distinguished where relevant. Collection cost is reported separately from
repeated diagnosis cost. Time summaries include the median and spread
across cases and repetitions, with timeout counts alongside completed-run
times. A fast failure must remain distinguishable from a fast solution.

Finally, incorrect and divergent outcomes are examined through their
recorded investigations. The analysis distinguishes failure to retrieve
available evidence, incorrect correlation, incorrect interpretation,
unsupported specificity and unsuccessful tool use. This explains where
a measured difference arises without assuming that the final score alone
identifies the mechanism.

\section{Remediation and Recovery Evaluation}
\label{sec:methodology:remediation}

Remediation is evaluated in the live environment because retained
telemetry cannot demonstrate a repair. Cases are selected for an
authorised action and a verifiable reset procedure before their
diagnostic results are known. This avoids measuring repair performance
only on cases an approach happens to solve. Each executed plan is linked
to its originating diagnosis and configuration.

Plan review checks the proposed target, parameters, expected effect,
preconditions, dry-run result and rollback against the independent case
record. The operator then approves or rejects that specific plan.
Reference-answer corrections are not inserted into the plan and counted
as agent success. A rejected or unsafe recommendation is retained in the
results without executing it merely to measure its harm.

Before each repair trial, the fault is re-established and its effect
verified. Approval waiting time and execution time are measured separately.
Recovery requires the affected operation to behave as expected through
the configured Cloudflare test hostname, together with checks of the relevant
resource state. A healthy probe alone is insufficient when the failed
operation remains broken. The expected responses, repeated verification
requests and observation interval are fixed in the case record. The
fault harness must not reset the fault during this check.

The applied outcome is successful, unsuccessful or harmful, following
Norlys AC-7 \parencite[AC-7]{norlys_operations_agents}. Success requires
verified recovery without an observed violation of the case's scope and
state checks. Harm means an observed prohibited change or additional
degradation within those checks, rather than a speculative consequence.
Command completion without recovery is unsuccessful. Results include
plan eligibility, rejection, approval, execution and recovery counts.
Recovery among executed plans and recovery among all repair-eligible cases
are reported separately, exposing failures before execution.

Rollback has a separate test. At least one controlled trial applies a
permitted change, invokes rollback and checks the saved configuration and
expected endpoint behaviour. Restoring the pre-repair state can restore
the original fault. Successful rollback therefore means returning to the
specified prior state, not necessarily keeping the endpoint healthy.
The harness subsequently resets the case to its healthy starting state.

Negative control tests exercise unapproved, rejected, modified and
out-of-scope plans, a viewer attempting approval and an unsupported action.
These tests use authorised test resources and check both the audit record
and absence of a state change. Fixed reference plans can test these
controls independently of diagnostic success, but are excluded from
agent repair-success results. The controls, dry run and rollback outcomes
are reported individually. Zero observed violations establishes behaviour
on these tests, not a general guarantee of safety.

\section{Validity and Reproducibility}
\label{sec:methodology:validity}

\noindent\textbf{Attributing differences.}
Shared evidence, models and total limits reduce competing explanations
for differences between the agents. Role instructions and communication
still differ, so conclusions concern the tested configurations. Fixed
evidence removes live collection variability. Live connector and
remediation tests address that separate limitation.

\noindent\textbf{Measuring the intended outcome.}
Independent reference records and separate component, cause and evidence
scores reduce the risk of rewarding plausible wording. A scripted fault
can still be easier to recognise than a naturally occurring incident.
Removing answer-revealing labels, retaining realistic symptoms and using
reserved cases limits this risk. Reviewer judgement remains necessary
for causal explanations, making the rubric and independent review part
of the evidence. The operator tests assess workflow and enforced
permissions, not general usability or human decision quality.

\noindent\textbf{Scope of the conclusions.}
The study uses one application path, one organisation and a deliberately
selected set of faults. Reserved cases reduce tuning to exact answers
but do not make the sample representative of production incidents.
Uneven fault coverage, unavailable security-control injections and the
low-traffic workload limit transfer to other systems. Repetitions expose
model variation within these cases but cannot compensate for missing
fault families. Interpretations remain bounded by the measured path,
model version and permitted actions.

Reproduction requires case descriptions, versions, prompts, rules, tool
contracts, the run schedule, scoring rubric and analysis code. Retained
observations and reference answers remain access-controlled under Norlys'
data policy. Public material can provide sanitised descriptions, schemas
and aggregate results. Each result identifies its case, evidence collection
and method version so that its calculation and supporting observations
can be checked.
